India has sparse $PM_{2.5}$ monitoring, causing unmonitored urban and suburban areas. Existing predictive models often have weak spatial transferability and tend to attenuate severe pollution peaks. To address this, we propose a two-stage framework for hourly $PM_{2.5}$ estimation. Stage 1 estimates a daily background $PM_{2.5}$ level, which is used to scale the hourly predictions generated by Stage 2 from hourly meteorology and static land-use features. The model incorporates physics-guided features, including inverse planetary boundary layer height ($1/\text{PBLH}$) and continuous wind vectors, and is optimized using a frequency-weighted Mean Absolute Error ($\mathcal{L}_{fMAE}$) alongside a variance penalty. We evaluate the framework using a zero-shot cross-city holdout on unseen cities (Lucknow and Patna) rather than random test samples. On the combined held-out-city evaluation, the full hybrid model improved $R^2$ to 0.6119 from the 0.4890 baseline established by Stage 1 alone. In a severe evaluated event, the model closely tracked the observed peak by generating a 6.725$\times$ hourly ratio, producing a predicted peak of approximately 846.2 $\mu g/m^3$ against an observed 830 $\mu g/m^3$. These results demonstrate that this low-compute framework has strong potential to support neighborhood-scale decision-support tools in unmonitored airsheds.
